\documentclass[11pt]{article}
\usepackage{amsmath,amssymb,amsthm}
\usepackage[margin=1.2in]{geometry}
\newtheorem{theorem}{Theorem}
\newtheorem{proposition}[theorem]{Proposition}
\newtheorem{corollary}[theorem]{Corollary}
\newtheorem{remark}[theorem]{Remark}
\DeclareMathOperator{\rad}{rad}
\title{Erd\H{o}s \#700: the capped carry ledger for
$\min_k\gcd\bigl(n,\binom nk\bigr)$}
\author{Samuel Lavery}
\date{Draft --- \today}
\begin{document}
\maketitle

\begin{abstract}
For $f(n)=\min_{1<k\le n/2}\gcd\bigl(n,\binom nk\bigr)$ we give the exact form of the
minimand as a product of per-prime carry counts \emph{capped} at $v_p(n)$, deduce
$f(n)\mid n$ and $f(n)\le\rad(n)$, and explain from the cap why the measured set
$\{n:f(n)>\sqrt n\}$ is squarefree-dominated.  \#700 is not proved.
\end{abstract}

\section{The capped ledger}
Let $\kappa_p(k,n-k)$ denote the number of carries in the base-$p$ addition $k+(n-k)$, so
that $v_p\binom nk=\kappa_p$ by Kummer.

\begin{proposition}\label{prop:capped}
For $n$ composite and $1<k\le n/2$,
\[
  \gcd\!\left(n,\binom nk\right)=\prod_{p^{a}\,\|\,n}p^{\min\left(a,\ \kappa_p(k,\,n-k)\right)},
  \qquad
  f(n)=\min_{1<k\le n/2}\ \prod_{p^{a}\,\|\,n}p^{\min(a,\,\kappa_p)} .
\]
In particular $f(n)\mid n$, and $f(n)\le\rad(n)$ when $n$ is squarefree.
\end{proposition}

\begin{proof}
The valuation of a gcd is the minimum of the valuations, and $v_p(n)=a$,
$v_p\binom nk=\kappa_p$.  Each exponent is at most $a$, so the product divides $n$.  If $n$
is squarefree every $a=1$ and each factor is $p^{0}$ or $p^{1}$.
\end{proof}

\emph{Verified}: $f(n)\mid n$ with no exception for all composite $n<300$.

\begin{remark}[why the cap matters]
This is the one member of the $\binom nk$ domination cluster whose rails \emph{saturate}:
rail $p$ contributes at most $v_p(n)$ however many carries occur.  Elsewhere in the cluster
(\#683, \#685, \#1094, \#1095) the relevant statistic is unbounded in the carry count, and
the analysis is driven by which rails carry at all; here it is driven by which rails can be
made \emph{not} to carry, simultaneously, by a single choice of $k$.
\end{remark}

\begin{corollary}\label{cor:upper}
For composite $n$, taking $k=p^{a}$ with $p^{a}\,\|\,n$ gives
$\gcd\bigl(n,\binom{n}{p^{a}}\bigr)=n/p^{a}$, whence $f(n)\le n/P(n)$ where $P(n)$ is the
largest prime factor, and therefore $f(n)\le(1+o(1))n/\log n$.
\end{corollary}

This is the bound recorded by Erd\H{o}s and Szekeres; it is included because
Proposition~\ref{prop:capped} makes it a one-line valuation computation.

\section{Prime powers, and the correct divisor}

Erd\H{o}s and Szekeres record $f(n)\le n/P(n)$ from the choice $k=p^{a}$ with $p^{a}\|n$.
That choice is legitimate only when $p^{a}\le n/2$, and tracking the restriction determines
the prime-power case completely.

\begin{theorem}\label{thm:pp}
Let $p$ be prime and $a\ge2$.  Then $f(p^{a})=p$.
\end{theorem}

\begin{proof}
Put $n=p^{a}$ and let $1<k\le n/2$.  Since $s_p(n)=1$, Legendre and Kummer give
\[
  v_p\!\binom nk=\kappa_p=\frac{s_p(k)+s_p(n-k)-1}{p-1},
\]
and $n$ has the single nonzero digit $1$ in position $a$, so any addition $k+(n-k)=p^{a}$
with $0<k<p^{a}$ carries at least once; hence $\kappa_p\ge1$ and $p\mid\binom nk$ for every
such $k$.  As $p$ is the only prime dividing $n$, this gives
$\gcd\bigl(n,\binom nk\bigr)=p^{\min(a,\kappa_p)}\ge p$, so $f(p^{a})\ge p$.

For the reverse, take $k=p^{a-1}$, which satisfies $1<k\le n/2$ because $a\ge2$ and
$p\ge2$.  Then $s_p(k)=1$ and $n-k=p^{a-1}(p-1)$ has the single digit $p-1$, so
$s_p(n-k)=p-1$ and
\[
  \kappa_p=\frac{1+(p-1)-1}{p-1}=1 .
\]
Hence $\gcd\bigl(n,\binom{n}{p^{a-1}}\bigr)=p^{\min(a,1)}=p$, so $f(p^{a})\le p$.
\end{proof}

\emph{Verified}: all prime powers $p^{a}\le2500$ with $p\le13$, $a\ge2$.

\begin{proposition}\label{prop:qbound}
Let $n$ be composite and let $q=p^{v_p(n)}$ be any prime power exactly dividing $n$ with
$q\le n/2$.  Then $f(n)\le n/q$.  Consequently
\[
  f(n)\ \le\ \frac{n}{Q'(n)},\qquad
  Q'(n):=\max\bigl\{p^{v_p(n)}\ :\ p^{v_p(n)}\le n/2\bigr\}.
\]
\end{proposition}

\begin{proof}
Take $k=q$ in the minimum defining $f(n)$; the constraint $1<k\le n/2$ holds by hypothesis.
Writing $q=p^{a}$ with $a=v_p(n)$, Kummer gives $\kappa_p(q,n-q)=a$, since adding $q$ to
$n-q$ in base $p$ carries out of each of the positions $0,\dots,a-1$ and no further, while
for $\ell\ne p$ one has $\kappa_\ell\ge0$.  Hence
$\gcd\bigl(n,\binom nq\bigr)=n/q$, and $f(n)$, being a minimum, is at most this.
\end{proof}

\emph{Verified}: $1095$ pairs $(n,q)$ with $n<600$.

\begin{corollary}\label{cor:necessity}
If $n$ is composite, is not the square of a prime, and satisfies $f(n)=n/P(n)$, then
$P(n)=Q'(n)$: the largest prime factor of $n$ is also its largest admissible prime power.
\end{corollary}

\begin{proof}
By Proposition~\ref{prop:qbound}, $f(n)\le n/Q'(n)\le n/P(n)$ whenever $P(n)\le n/2$, which
holds for composite $n$ other than $n=p^{2}$.  Equality $f(n)=n/P(n)$ therefore forces
$Q'(n)\le P(n)$, and $Q'(n)\ge P(n)$ always, since $P(n)$ itself is an admissible prime
power in that range.
\end{proof}

\emph{Verified}: the exceptions to $P=Q'$ among $n<600$ with $f(n)=n/P(n)$ are exactly the
prime squares $4,9,25,49,121,169,289,361,529$, and no others.

\section{Measurement}
Census over composite $n<600$:
\begin{itemize}
\item $f(n)=n/P(n)$ for $307$ values, beginning
      $4,6,9,10,14,15,20,21,22,25,26,28,30,33$ --- consistent with the known cases
      (products of two primes, and $n=30$).
\item $f(n)>\sqrt n$ for $30$ values: $30,70,84,105,132,140,154,165,168,182,\dots$ with
      $f=6,10,12,15,12,14,14,15,21,14$.  The count is still growing at the top of the range.
\end{itemize}
The second list is squarefree-dominated, which Proposition~\ref{prop:capped} explains: for
$f(n)$ to be large every rail must reach its cap for \emph{every} $k$, and a squarefree $n$
with several small prime factors is the configuration in which all caps are $1$ and can be
attained together.

\begin{remark}[register]
\textbf{Proved}: Proposition~\ref{prop:capped} and Corollary~\ref{cor:upper}, elementary from
Legendre and Kummer; neither is claimed new (Corollary~\ref{cor:upper} is Erd\H{o}s--Szekeres).
Theorem~\ref{thm:pp} determines $f$ on prime powers, $f(p^{a})=p$ for $a\ge2$;
Proposition~\ref{prop:qbound} replaces $P(n)$ by the largest \emph{admissible prime power}
$Q'(n)$; and Corollary~\ref{cor:necessity} gives a necessary condition,
$P(n)=Q'(n)$, for the equality $f(n)=n/P(n)$ that Erd\H{o}s asks to characterise, with the
prime squares as the exact exceptions.
\textbf{Measured}: the two censuses.  \textbf{Not proved}: the characterisation of
$\{n:f(n)=n/P(n)\}$, the infinitude of $\{n:f(n)>\sqrt n\}$, or the bound
$f(n)\ll_A n/(\log n)^{A}$.  \textbf{Falsifier}: $\{n:f(n)>\sqrt n\}$ terminating at larger
$n$ would contradict the apparent growth.  \textbf{Next}: census to $n=10^{5}$ and test
whether that set has positive density; and characterise $\{f(n)=n/P(n)\}$ by $\Omega(n)$ and
squarefreeness, where the data suggests a clean statement is available.
\end{remark}
\end{document}
